\begin{definition}[Truncated norm map at stationary depth]
\label{mod:parameters-normpolynomial}
Let $K/F$ be a finite free valuative extension of nonarchimedean local
fields.  A stationary conductor decomposition is certified data
\[
 m=2d+\varepsilon>1,\qquad 0\leq\varepsilon\leq1.
\]
It supplies
\[
 r=d+\varepsilon>0,\qquad r\leq m,\qquad m\leq2r,
 \qquad m-r=d.
\]
For decompositions
\[
 m_K=2d_K+\varepsilon_K,\qquad
 m_F=2d_F+\varepsilon_F,
\]
put $r_K=d_K+\varepsilon_K$ and $r_F=d_F+\varepsilon_F$.

For natural-number depths $a,b$, let
\[
 \operatorname{NormIncl}_{K/F}(a,b)
\]
mean the exact filtration inclusion
\[
 N_{K/F}(U_K^a)\subseteq U_F^b.
\]
A norm-polynomial precision certificate retains separately the three
inclusions
\begin{equation}\label{eq:lean-norm-polynomial-three-inclusions}
 \begin{aligned}
 N(U_K^{r_K})&\subseteq U_F^{r_F},\\
 N(U_K^{m_K})&\subseteq U_F^{m_F},\\
 N(U_K^{d_K})&\subseteq U_F^{d_F}.
 \end{aligned}
\end{equation}
The first gives target membership, the second gives independence modulo
the conductor layer, and the third is the additional stationary-numerator
ambiguity inclusion.  These are distinct fields of
\texttt{NormPolynomialPrecision}.

More generally, suppose
\[
 0<a_K\leq b_K\leq2a_K,\qquad
 0<a_F\leq b_F\leq2a_F,
\]
and assume
\[
 N(U_K^{a_K})\subseteq U_F^{a_F},
 \qquad
 N(U_K^{b_K})\subseteq U_F^{b_F}.
\]
Norm first descends to the honest unit-filtration quotients and, through
the positive-depth unit--lattice equivalences, induces the
representative-free additive homomorphism
\[
 \operatorname{trNorm}_{K/F}:
 \mathfrak p_K^{a_K}/\mathfrak p_K^{b_K}
 \longrightarrow
 \mathfrak p_F^{a_F}/\mathfrak p_F^{b_F}.
\]
On every supplied numerator it is
\[
 [x]\longmapsto[N_{K/F}(1+x)-1].
\]
Changing $x$ modulo $\mathfrak p_K^{b_K}$ changes $1+x$ by a ratio in
$U_K^{b_K}$, so the second inclusion makes the output independent of the
source representative.  Exact additivity follows from
\[
 (1+x)(1+y)=1+x+y+xy.
\]
Indeed,
$xy\in\mathfrak p_K^{2a_K}\subseteq\mathfrak p_K^{b_K}$, while the
product of the two output displacements lies in
$\mathfrak p_F^{2a_F}\subseteq\mathfrak p_F^{b_F}$.

The principal declaration \texttt{normPolynomial} is the stationary
specialization
\[
 P_{K/F}:
 \mathfrak p_K^{r_K}/\mathfrak p_K^{m_K}
 \longrightarrow
 \mathfrak p_F^{r_F}/\mathfrak p_F^{m_F}.
\]
Both sides are the actual finite lattice quotients used by Lamprecht
duality.  Evaluation, zero, addition, and source-representative
independence are exported separately.

For $n=[K:F]$, retain the exact polynomial value
\begin{equation}\label{eq:lean-exact-norm-polynomial}
 P_{K/F}(x)=N_{K/F}(1+x)-1
           =\sum_{j=1}^{n}s_j(x),
\end{equation}
where
\[
 s_1(x)=\operatorname{Tr}_{K/F}(x),
 \qquad s_n(x)=N_{K/F}(x).
\]
Equation~\ref{eq:lean-exact-norm-polynomial} also holds after passage to
the target quotient.  Define the complete nonlinear correction
\[
 C_{K/F}(x)=\sum_{j=2}^{n}s_j(x);
\]
then the exact identity is
\[
 P_{K/F}(x)=\operatorname{Tr}_{K/F}(x)+C_{K/F}(x).
\]
The trace descends to the lattice quotients only after separately
certifying
\[
 \operatorname{Tr}(\mathfrak p_K^{r_K})
   \subseteq\mathfrak p_F^{r_F},
 \qquad
 \operatorname{Tr}(\mathfrak p_K^{m_K})
   \subseteq\mathfrak p_F^{m_F}.
\]

The strong wild estimate from
Lemma~\ref{lem:wild-symmetric-bound} supplies vanishing only for the
intermediate degrees $2\leq j<n$.  In a totally ramified extension the
terminal bound is supplied separately by
$v_F(s_n(x))=v_K(x)$.  Consequently intermediate vanishing alone gives
the critical formula
\[
 P_{K/F}(x)\equiv
 \operatorname{Tr}_{K/F}(x)+N_{K/F}(x)
 \pmod{\mathfrak p_F^{m_F}}.
\]
The trace-only identity
\[
 P_{K/F}=\operatorname{Tr}_{K/F}
 \pmod{\mathfrak p_F^{m_F}}
\]
is proved only when every intermediate coefficient and, independently,
the terminal norm coefficient belong to $\mathfrak p_F^{m_F}$.  Thus the
critical norm polynomial is never identified with its trace term outside
the range justified by the symmetric bounds.

The predicate \texttt{NormPolynomialElementaryBoundsAt} records that
every term in the exact expansion has a prescribed target valuation.
Its degree-one member gives the trace-lattice inclusion, while bounds on
all terms give the corresponding norm-filtration inclusion.  Bounds at
$(r_K,r_F)$, $(m_K,m_F)$, and $(d_K,d_F)$ therefore construct the full
three-part certificate in
\eqref{eq:lean-norm-polynomial-three-inclusions}.

The ambiguity inclusion induces the honest multiplicative quotient map
\[
 U_K^0/U_K^{d_K}\longrightarrow U_F^0/U_F^{d_F},
 \qquad [u]\longmapsto[N_{K/F}(u)].
\]
This is \texttt{normPolynomialStationaryNormClass}.  In particular, the
norm of a stationary unit class is independent of its unit representative
at precisely the extra ambiguity depth.  No norm map on arbitrary
representatives of the additive quotient is asserted.

Both the unit-quotient norm and the additive truncated norm commute with
replacing stronger denominator depths by shallower ones.  For nested
source and target denominator depths, the square formed by the canonical
quotient projections and the corresponding truncated norm maps commutes.
The construction preserves supplied numerators but makes no quotient
representative choice.

Finally, retain the denominator data required by the later adjoint.  For
$\gamma_F\in F^\times$ and $\gamma_K\in K^\times$, define
\[
 \rho_{\gamma_F,\gamma_K}
   =\frac{\gamma_K}{\iota_{F,K}(\gamma_F)}\in K^\times.
\]
Its nonvanishing is explicit.  If
\[
 \operatorname{ord}_K(\rho_{\gamma_F,\gamma_K})=0,
 \qquad d_K\leq e(K/F)d_F,
\]
then the ramification scaling
\[
 \operatorname{ord}_K(\iota_{F,K}(x))
   =e(K/F)\operatorname{ord}_F(x)
\]
makes
\[
 \mathcal O_F/\mathfrak p_F^{d_F}
 \longrightarrow
 \mathcal O_K/\mathfrak p_K^{d_K},
 \qquad
 [x]\longmapsto
 [\rho_{\gamma_F,\gamma_K}\,\iota_{F,K}(x)]
\]
a representative-free additive homomorphism.  For the common denominator
$\gamma_K=\iota_{F,K}(\gamma_F)$ it is the ordinary base-field inclusion.
The exact identity required by the trace adjoint is
\[
 \operatorname{Tr}_{K/F}\!\left(
   \frac{\rho_{\gamma_F,\gamma_K}\,
     \iota_{F,K}(x)y}{\gamma_K}\right)
 =
 \frac{x\,\operatorname{Tr}_{K/F}(y)}{\gamma_F}.
\]
Thus the truncated norm itself is denominator-independent, while every
denominator-sensitive scaling used by its later adjoint retains the
specified pair $(\gamma_F,\gamma_K)$.

\LeanFile{LanglandsFirstMainLemma/Parameters/NormPolynomial.lean}
\lean{LanglandsFirstMainLemma.IsStationaryConductorDecomposition}
\lean{LanglandsFirstMainLemma.IsStationaryConductorDecomposition.variableDepth_pos}
\lean{LanglandsFirstMainLemma.IsStationaryConductorDecomposition.floorDepth_pos}
\lean{LanglandsFirstMainLemma.IsStationaryConductorDecomposition.variableDepth_le_conductor}
\lean{LanglandsFirstMainLemma.IsStationaryConductorDecomposition.conductor_le_two_variableDepth}
\lean{LanglandsFirstMainLemma.IsStationaryConductorDecomposition.conductor_sub_variableDepth}
\lean{LanglandsFirstMainLemma.NormUnitFiltrationInclusion}
\lean{LanglandsFirstMainLemma.NormPolynomialPrecision}
\lean{LanglandsFirstMainLemma.normUnitFiltrationHom}
\lean{LanglandsFirstMainLemma.coe_normUnitFiltrationHom}
\lean{LanglandsFirstMainLemma.normUnitFiltrationQuotientHom}
\lean{LanglandsFirstMainLemma.normUnitFiltrationQuotientHom_mk}
\lean{LanglandsFirstMainLemma.truncatedNormPolynomial}
\lean{LanglandsFirstMainLemma.truncatedNormPolynomial_mk}
\lean{LanglandsFirstMainLemma.normPolynomial}
\lean{LanglandsFirstMainLemma.normPolynomialValue}
\lean{LanglandsFirstMainLemma.normPolynomialValue_eq_sum_elementarySymmetric}
\lean{LanglandsFirstMainLemma.normPolynomialRepresentative}
\lean{LanglandsFirstMainLemma.coe_normPolynomialRepresentative}
\lean{LanglandsFirstMainLemma.normPolynomial_mk}
\lean{LanglandsFirstMainLemma.normPolynomial_zero}
\lean{LanglandsFirstMainLemma.normPolynomial_add}
\lean{LanglandsFirstMainLemma.normPolynomial_representative_independent}
\lean{LanglandsFirstMainLemma.normPolynomial_mk_eq_sum_elementarySymmetric}
\lean{LanglandsFirstMainLemma.TraceLatticeInclusion}
\lean{LanglandsFirstMainLemma.traceLatticeHom}
\lean{LanglandsFirstMainLemma.coe_traceLatticeHom}
\lean{LanglandsFirstMainLemma.traceLatticeQuotientHom}
\lean{LanglandsFirstMainLemma.traceLatticeQuotientHom_mk}
\lean{LanglandsFirstMainLemma.sum_range_pred_shift_eq_sum_Icc}
\lean{LanglandsFirstMainLemma.normPolynomialHigherCorrectionValue}
\lean{LanglandsFirstMainLemma.normPolynomialValue_eq_trace_add_higherCorrection}
\lean{LanglandsFirstMainLemma.NormPolynomialIntermediateTermsVanishAt}
\lean{LanglandsFirstMainLemma.NormPolynomialTerminalTermVanishAt}
\lean{LanglandsFirstMainLemma.wild_normPolynomialIntermediateTermsVanishAt}
\lean{LanglandsFirstMainLemma.totallyRamified_normPolynomialTerminalTermVanishAt}
\lean{LanglandsFirstMainLemma.normPolynomialValue_sub_trace_mem}
\lean{LanglandsFirstMainLemma.normPolynomialValue_sub_trace_add_norm_mem}
\lean{LanglandsFirstMainLemma.normPolynomial_mk_eq_trace_add_norm}
\lean{LanglandsFirstMainLemma.normPolynomialTrace}
\lean{LanglandsFirstMainLemma.normPolynomialTrace_mk}
\lean{LanglandsFirstMainLemma.normPolynomial_eq_trace}
\lean{LanglandsFirstMainLemma.NormPolynomialElementaryBoundsAt}
\lean{LanglandsFirstMainLemma.traceLatticeInclusion_of_elementarySymmetricBounds}
\lean{LanglandsFirstMainLemma.normUnitFiltrationInclusion_of_elementarySymmetricBounds}
\lean{LanglandsFirstMainLemma.NormPolynomialPrecision.ofElementarySymmetricBounds}
\lean{LanglandsFirstMainLemma.normUnitFiltrationInclusion_zero}
\lean{LanglandsFirstMainLemma.stationaryNormClass}
\lean{LanglandsFirstMainLemma.stationaryNormClass_mk}
\lean{LanglandsFirstMainLemma.stationaryNormClass_representative_independent}
\lean{LanglandsFirstMainLemma.normPolynomialStationaryNormClass}
\lean{LanglandsFirstMainLemma.normUnitFiltrationQuotientHom_projection}
\lean{LanglandsFirstMainLemma.truncatedNormPolynomial_projection}
\lean{LanglandsFirstMainLemma.normPolynomialDenominatorRatio}
\lean{LanglandsFirstMainLemma.normPolynomialDenominatorRatio_ne_zero}
\lean{LanglandsFirstMainLemma.denominatorScaledRepresentative}
\lean{LanglandsFirstMainLemma.denominatorScaledAlgebraMap}
\lean{LanglandsFirstMainLemma.denominatorScaledAlgebraMap_mk}
\lean{LanglandsFirstMainLemma.denominatorScaledAlgebraMap_common_mk}
\lean{LanglandsFirstMainLemma.trace_denominatorScaled_div}
\leanok
\SourceText{Lemma~\ref{lem:stationary-depth-norm},
equation~\ref{eq:norm-expansion}, and the truncated norm construction in
Lemma~\ref{lem:truncated-norm-adjoint}.}
\uses{mod:localfield-extension,mod:ramification-symmetricbounds,mod:lamprecht-stationaryclass}
\FormalizationContract
The source and target are certified finite lattice quotients.  The
principal norm, projection, stationary-unit, and denominator-scaling maps
consume quotient classes; representative-level declarations only evaluate
those maps on a supplied numerator and never choose a quotient
representative.  The variable, conductor, and ambiguity norm
inclusions remain distinct.  Exact additivity records both half-depth
inequalities.  Trace compatibility requires separate numerator and
denominator trace inclusions together with vanishing of every intermediate
symmetric term and the terminal norm term.  Denominator-ratio scaling
retains its chosen pair, exact local-unit order, and ramification-scaled
depth.
\end{definition}
